\chapter{DH 参数的几何定义}

\section{连杆与关节的编号}

在 DH 参数法中，约定如下编号规则：

\begin{figure}[htbp]
	\centering
	\begin{tikzpicture}[scale=1.1]
		% 基座
		\fill[pattern=north east lines, pattern color=gray!50] (-0.5,0) rectangle (0.5,0.5);
		\node at (0,0.25) {基座 0};
		
		% 关节1轴
		\draw[ultra thick, blue] (0,0.5) -- (0,1.5);
		\node[left] at (0,1.0) {轴 $i-1$};
		
		% 连杆 i-1
		\draw[very thick] (0,0.5) -- (0,1.5);
		
		% 关节 i-1
		\draw[fill=orange!30] (-0.2,0.5) rectangle (0.2,0.7);
		\node[left] at (-0.2,0.6) {关节 $i-1$};
		
		% 连杆 i
		\draw[very thick, red!60!black] (0,1.5) -- (3.5,2.5);
		\node[red!60!black] at (1.5,2.0) {连杆 $i$};
		
		% 关节 i
		\draw[fill=orange!30, rotate around={30:(3.5,2.5)}] (3.3,2.3) rectangle (3.7,2.7);
		\node[above] at (3.8,2.6) {关节 $i$};
		
		% 关节i轴
		\draw[ultra thick, blue, rotate around={30:(3.5,2.5)}] (3.5,2.0) -- (3.5,3.5);
		\node[right] at (4.2,2.9) {轴 $i$};
		
		% 连杆 i+1
		\draw[very thick, red!60!black] (3.5,2.5) -- (5.5,3.5);
		
		% 关节 i+1
		\draw[fill=orange!30, rotate around={20:(5.5,3.5)}] (5.3,3.3) rectangle (5.7,3.7);
		\node[right] at (5.8,3.6) {关节 $i+1$};
		
		% 公垂线
		\draw[dashed, green!50!black, thick] (0,1.5) -- (3.5,2.5);
		\node[green!50!black] at (2.2,2.9) {$a_{i-1}$（公垂线）};
	\end{tikzpicture}
	\caption{连杆与关节的编号约定}
\end{figure}

\textbf{编号规则：}
\begin{itemize}
	\item 基座编号为 0，视为\textbf{虚拟连杆 0}。
	\item 从基座向外，依次编号：\textbf{关节 1、连杆 1、关节 2、连杆 2、……、关节 n、连杆 n}。
	\item \textbf{连杆 $i$} 位于关节 $i$ 和关节 $i+1$ 之间。
	\item 关节 $i$ 连接连杆 $i-1$ 和连杆 $i$；关节 $i$ 运动时，连杆 $i$ 随之运动。
	\item 末端执行器（工具）编号为 $n+1$，视为\textbf{虚拟连杆 $n+1$}。
\end{itemize}

\section{四个 DH 参数的几何意义}

对于每一对相邻的关节轴（轴 $i-1$ 与轴 $i$），四个 DH 参数完整刻画了它们的空间几何关系。

\begin{mydefinition}{四个 DH 参数（标准 DH）}
	设关节轴 $i-1$ 和关节轴 $i$ 为空间中两条直线，定义：
	\begin{enumerate}
		\item \textbf{连杆长度} $a_{i-1}$：沿 $x_{i-1}$ 轴方向，轴 $i-1$ 与轴 $i$ 之间的公垂线长度。当两轴平行时，有无数条等长公垂线；当两轴相交时，$a_{i-1}=0$。\\
		\item \textbf{连杆扭角} $\alpha_{i-1}$：绕 $x_{i-1}$ 轴，从轴 $i-1$ 方向旋转到轴 $i$ 方向的角度（右旋为正）。\\
		\item \textbf{连杆偏距} $d_i$：沿 $z_{i-1}$ 轴方向，从 $x_{i-1}$ 到 $x_i$ 的有向距离。\\
		\item \textbf{关节转角} $\theta_i$：绕 $z_{i-1}$ 轴，从 $x_{i-1}$ 方向旋转到 $x_i$ 方向的角度（右旋为正）。
	\end{enumerate}
\end{mydefinition}

\begin{figure}[htbp]
	\centering
	\begin{tikzpicture}[scale=1.3]
		% 轴 i-1 (垂直向上)
		\draw[->, ultra thick, blue] (0,-0.5) -- (0,3.5) node[above] {轴 $i-1$};
		\node[left] at (0,1.5) {$z_{i-1}$};
		
		% 轴 i (倾斜)
		\draw[->, ultra thick, blue, rotate around={25:(3,0.5)}] (3,-0.5) -- (3,3.5) node[above] {轴 $i$};
		\node at (2.5,2.2) {$z_{i}$};
		
		% 公垂线 = a_{i-1}
		\draw[->, very thick, red!70!black] (0,1) -- node[midway, above] {$a_{i-1}$} (3,1);
		\node[below] at (1.5,1) {$x_{i-1}$};
		
		% x_{i-1}轴延伸
		\draw[->, thick, red!70!black] (0,1) -- (-0.8,1);
		
		% 原点 O_{i-1}
		\filldraw (0,1) circle (1.5pt) node[below left] {$O_{i-1}$};
		
		% 交点 O_i
		\filldraw (3,1) circle (1.5pt) node[below right] {$O_{i}$};
		
		% x_i 轴（从O_i沿公垂线方向）
		\draw[->, thick, red!70!black] (3,1) -- (4,1) node[right] {$x_i$};
		
		% 标注 alpha_{i-1}
		\draw[->, thick, green!50!black] (2.6,1.5) arc[start angle=90, end angle=115, radius=1.2];
		\node[green!50!black] at (2.3,1.8) {$\alpha_{i-1}$};
		
		% z_{i-1}在O_{i-1}处
		\draw[->, thick, blue!50!black] (0,1) -- (0,2.2);
		
		% z_i在O_i处
		\begin{scope}[shift={(3,1)}, rotate=25]
			\draw[->, thick, blue!50!black] (0,0) -- (0,1.5);
		\end{scope}
		
		% 标注 d_i (偏距，沿z_{i-1}方向)
		\draw[<->, dashed, purple] (-0.5,0.5) -- (-0.5,1) node[midway, left] {$d_i$};
		
		% 标注 theta_i (绕z_{i-1}的转角)
		\draw[->] (0.35,1.25) arc[start angle=0, end angle=20, radius=0.8];
		\node at (0.8,1.6) {$\theta_i$};
		
	\end{tikzpicture}
	\caption{标准 DH 参数的几何意义（空间示意）}
\end{figure}

\subsection{参数的进一步解释}

\begin{enumerate}
	\item \textbf{$a_{i-1}$（连杆长度，Link Length）}：沿 $x_{i-1}$ 轴方向，从 $z_{i-1}$ 到 $z_i$ 的\textbf{公垂线长度}，即两相邻关节轴线之间的最短距离。恒为非负且永远为常数。若两轴相交，$a_{i-1}=0$；若两轴平行，可在无穷多条等长公垂线中任选一条，通常选与前一公垂线相交的那一条以保证 $d_i$ 的确定性。
	
	\item \textbf{$\alpha_{i-1}$（连杆扭角，Link Twist）}：绕 $x_{i-1}$ 轴旋转、将 $z_{i-1}$ 转向 $z_i$ 所需的角度。描述了两关节轴线在空间中的夹角，恒为常数，范围通常取 $[-\pi, \pi]$。右手定则：右手拇指指向 $x_{i-1}$ 正向，四指弯曲方向为 $\alpha_{i-1}$ 的正方向。两轴平行同向时 $\alpha_{i-1}=0$，平行反向时 $\alpha_{i-1}=\pi$。
	
	\item \textbf{$d_i$（连杆偏距，Link Offset）}：沿 $z_{i-1}$ 轴方向，从 $x_{i-1}$ 到 $x_i$ 的轴向距离，描述相邻两根公垂线在关节轴上的偏移量。沿 $z_{i-1}$ 正向为正。对于\textbf{移动关节}，$d_i$ 是关节变量；对于\textbf{旋转关节}，$d_i$ 是常数。
	
	\item \textbf{$\theta_i$（关节转角，Joint Angle）}：绕 $z_{i-1}$ 轴旋转、将 $x_{i-1}$ 转向 $x_i$ 所需的角度。对于\textbf{旋转关节}，$\theta_i$ 是关节变量；对于\textbf{移动关节}，$\theta_i$ 是常数。
\end{enumerate}

\begin{myTheorem}{DH 参数的变量性}
	\begin{center}
		\begin{tabular}{ccc}
			\toprule
			关节类型 & 关节变量 & 常量参数 \\
			\midrule
			旋转关节 (R) & $\theta_i$ & $a_{i-1}, \alpha_{i-1}, d_i$ \\
			移动关节 (P) & $d_i$     & $a_{i-1}, \alpha_{i-1}, \theta_i$ \\
			\bottomrule
		\end{tabular}
	\end{center}
\end{myTheorem}

\textbf{注意：}无论哪种关节类型，$a_{i-1}$ 和 $\alpha_{i-1}$ 始终是常量（由机械结构决定），它们是对连杆本身的描述。而 $\theta_i$ 和 $d_i$ 中恰好有一个描述关节运动。

\subsection{几何直觉：两步对准}

将相邻连杆的连接理解为\textbf{两步空间对准}过程：

\begin{enumerate}
	\item \textbf{空间对齐轴线}（由 $a_{i-1}$、$\alpha_{i-1}$ 完成）：改变 $z$ 轴的位置与方向，使轴 $i-1$（$z_{i-1}$）与轴 $i$（$z_i$）重合。这一步由连杆本身的几何形状决定，对应 $T_x(a_{i-1})$ 和 $R_x(\alpha_{i-1})$ 两个变换。
	\item \textbf{沿轴线微调}（由 $d_i$、$\theta_i$ 完成）：在两 $z$ 轴已重合的前提下，沿共享的 $z$ 轴进行平移和旋转，使 $x_{i-1}$ 与 $x_i$ 对齐。这一步描述关节运动，对应 $R_z(\theta_i)$ 和 $T_z(d_i)$ 两个变换。
\end{enumerate}

对应的变换顺序（标准DH）为：

\begin{myformula}
	{}^{i-1}_{i}\mathbf{T} = R_z(\theta_i) \cdot T_z(d_i) \cdot T_x(a_{i-1}) \cdot R_x(\alpha_{i-1})
\end{myformula}

其物理含义为：先在关节 $i-1$ 处绕 $z_{i-1}$ 旋转 $\theta_i$，沿 $z_{i-1}$ 平移 $d_i$，再沿新 $x$ 轴平移 $a_{i-1}$，最后绕新 $x$ 轴旋转 $\alpha_{i-1}$，即可到达坐标系 $\{i\}$。

\subsection{参数确定口诀}

\begin{center}
	\fbox{\bfseries 绕 $z$ 转 $\theta$，沿 $z$ 移 $d$，沿 $x$ 移 $a$，绕 $x$ 转 $\alpha$}
\end{center}

按此顺序逐次递推坐标系即可得到 ${}^{i-1}_{i}\mathbf{T}$。关节类型速判：
\begin{itemize}
	\item 旋转关节：$\theta_i$ 为变量，$d_i$ 为常数
	\item 移动关节：$d_i$ 为变量，$\theta_i$ 为常数
	\item $a_{i-1}$、$\alpha_{i-1}$ \textbf{无论如何都是常数}
\end{itemize}

\section{连杆坐标系的建立规则}

DH 参数法的关键步骤是\textbf{为每个连杆固连一个坐标系}。坐标系的放置必须遵循统一规则。

\begin{mydefinition}{连杆坐标系 $\{i\}$ 的建立（标准 DH）}
	坐标系 $\{i\}$ 固连于连杆 $i$，按如下规则定义：
	\begin{enumerate}
		\item \textbf{原点 $O_i$}：关节轴 $i$ 与公垂线 $a_i$ 的交点（即轴 $i$ 和轴 $i+1$ 公垂线与轴 $i$ 的交点）。
		\item \textbf{$z_i$ 轴}：沿关节轴 $i$ 的方向（指向任意，但建议一致朝向末端）。
		\item \textbf{$x_i$ 轴}：沿从轴 $i$ 到轴 $i+1$ 的公垂线方向，从关节 $i$ 指向关节 $i+1$。
		\item \textbf{$y_i$ 轴}：由右手定则确定：$\hat{\mathbf{y}}_i = \hat{\mathbf{z}}_i \times \hat{\mathbf{x}}_i$。
	\end{enumerate}
\end{mydefinition}

\begin{figure}[htbp]
	\centering
	\begin{tikzpicture}[scale=1.2]
		% 画两个关节轴
		% 轴 i (垂直)
		\draw[->, ultra thick, blue] (0,-0.3) -- (0,4) node[above] {轴 $i$};
		\node[left] at (0,2.5) {$z_i$};
		
		% 轴 i+1 (倾斜)
		\draw[->, ultra thick, blue] (3.5,-0.3) -- (3.5,4) node[above] {轴 $i+1$};
		
		% 轴i略倾斜展示（稍微倾斜以显示三维感）
		% 实际上保持垂直以展示二维投影
		
		% 公垂线 a_i
		\draw[->, very thick, red!70!black] (0,1.5) -- node[midway, above] {$a_i$} (3.5,1.5);
		\node[below] at (1.75,1.5) {$x_i$};
		
		% O_i
		\filldraw (0,1.5) circle (1.5pt) node[left] {$O_i$};
		
		% 坐标系{i}
		\draw[->, thick, red!70!black] (0,1.5) -- (1.2,1.5) node[below] {$x_i$};
		\draw[->, thick, green!50!black] (0,1.5) -- (0,2.5) node[left] {$z_i$};
		\draw[->, thick, orange] (0,1.5) -- (-0.5,2.0) node[left] {$y_i$};
		
		% O_{i+1} （如果a_{i+1}定义的话）
		\filldraw (3.5,1.5) circle (1.5pt) node[right] {$O_{i-1}$ 型};
		
		% 轴 i 方向在O_i处的箭头
		\draw[->, thick, blue!50!black] (0,1.5) -- (0,2.5);
		
		% 轴 i+1方向
		\draw[->, thick, blue!50!black] (3.5,1.5) -- (3.5,2.5);
		\node[right] at (3.5,2.5) {$z_{i+1}$};
		
		% 标注alpha_i
		\draw[->, thick, green!50!black] (3.2,2.0) arc[start angle=90, end angle=115, radius=0.8];
		\node[green!50!black] at (2.8,2.3) {$\alpha_i$};
	\end{tikzpicture}
	\caption{连杆坐标系 $\{i\}$ 的建立：$z_i$ 沿关节轴 $i$，$x_i$ 沿公垂线指向关节 $i+1$}
\end{figure}

\subsection{基坐标系 $\{0\}$ 和末端坐标系 $\{n\}$ 的处理}

\begin{itemize}
	\item \textbf{基坐标系 $\{0\}$}：可任意放置，习惯上使 $z_0$ 沿关节轴 1 的方向。当关节变量 $\theta_1 = 0$（旋转关节）或 $d_1 = 0$（移动关节）时，使 $\{0\}$ 与 $\{1\}$ 重合。这样可令某些参数取零，简化表达式。
	
	\item \textbf{末端坐标系 $\{n\}$}：固连于末端连杆 $n$。通常将原点 $O_n$ 放在夹爪中心或工具参考点，$z_n$ 沿接近方向（approach），$x_n$ 沿法向或滑动方向。确保参数少、物理意义清晰即可。
\end{itemize}

\section{DH 参数表的建立}

将每个连杆-关节对的四个参数整理成表格，是 DH 参数法的标准输出形式：

\begin{center}
	\begin{tabular}{cccccc}
		\toprule
		连杆 $i$ & $a_{i-1}$ & $\alpha_{i-1}$ & $d_i$ & $\theta_i$ & 关节类型 \\
		\midrule
		1 & $a_0$   & $\alpha_0$   & $d_1$   & $\theta_1$ & R/P \\
		2 & $a_1$   & $\alpha_1$   & $d_2$   & $\theta_2$ & R/P \\
		$\vdots$ & $\vdots$ & $\vdots$ & $\vdots$ & $\vdots$ & $\vdots$ \\
		$n$ & $a_{n-1}$ & $\alpha_{n-1}$ & $d_n$ & $\theta_n$ & R/P \\
		\bottomrule
	\end{tabular}
\end{center}

\textbf{注：}上表中 $a_{i-1}$ 和 $\alpha_{i-1}$ 描述连杆 $i-1$ 本身的形状，$d_i$ 和 $\theta_i$ 描述连杆 $i-1$ 与连杆 $i$ 之间的连接关系。这种"错位"的编号方式在标准 DH 中常常令人困惑，但也正是因为这样，四个参数恰好对应从 $\{i-1\}$ 到 $\{i\}$ 的四个变换步骤。

\section{完整的参数赋值步骤}

\begin{myTheorem}{DH 参数赋值的操作流程}
	\begin{enumerate}
		\item \textbf{确定所有关节轴}：在机械图中为每个关节画出无限延伸的轴线（空间中两条直线）。
		\item \textbf{标注公垂线}：对每一对相邻关节轴，确定（或选取）公垂线。
		\item \textbf{建立坐标系}：依次在公垂线与关节轴的交点处放置坐标系原点，依规则定轴。
		\item \textbf{读参数}：从坐标系间的几何关系中读出 $a_{i-1}, \alpha_{i-1}, d_i, \theta_i$。
		\item \textbf{填表}：将结果填入 DH 参数表。
	\end{enumerate}
\end{myTheorem}

\subsection{特例处理}

\begin{itemize}
	\item \textbf{两轴平行}：有无穷多条公垂线，通常选择与上一公垂线相交的那一条（以保证 $d_i$ 的确定性）。
	\item \textbf{两轴相交}：$a_{i-1}=0$，坐标系原点取在交点处。$x_{i-1}$ 取为与 $z_{i-1}$ 和 $z_i$ 均垂直的方向。
	\item \textbf{第一轴}：坐标系 $\{0\}$ 可自由放置，一般选 $O_0$ 在关节轴 1 上，$z_0$ 沿关节轴 1。
\end{itemize}
